\chapter{Vươn Lên Từ Đáy Sâu (Rising from the Depths)}
\begin{center}
  \textit{\color{truyengold}``Không có thất bại nào là thất bại thực sự --- trừ khi ta từ bỏ hoàn toàn.''}\\[4pt]
  \textit{(No failure is a true failure --- unless you give up completely.)}\\[4pt]
  \small\color{truyenblue}--- Elbert Hubbard ---
\end{center}
\ngancach

\section{Hồ Chí Minh Trong Nhà Tù Tưởng Giới Thạch}
\begin{truyen}{Hồ Chí Minh Trong Nhà Tù Tưởng Giới Thạch}{Hồ Chí Minh --- Trung Quốc, 1942--1943}
\chuhoa{N}{ăm} 1942, Hồ Chí Minh bị chính quyền Tưởng Giới Thạch bắt giam không lý do, bị giải qua 18 nhà lao ở tỉnh Quảng Tây trong hơn một năm.\\[4pt]
\textit{(In 1942, Ho Chi Minh was imprisoned without charge by the Chiang Kai-shek government, transferred through 18 prisons in Guangxi province over more than a year.)}

Bị xích tay, chân mang gông, đối mặt với cái chết có thể đến bất kỳ lúc nào --- nhưng ông không ngừng sáng tác.\\[4pt]
\textit{(Handcuffed, feet in shackles, facing death that could come at any moment --- but he never stopped composing.)}

Ông viết \textit{Nhật Ký Trong Tù} --- tập thơ chữ Hán gồm 134 bài --- ngay trong lao tù.\\[4pt]
\textit{(He wrote \textit{Prison Diary} --- a collection of 134 Chinese-character poems --- right in prison.)}

Bài thơ ``Tự Khuyên Mình'' ông viết: ``Ví không có cảnh đông tàn, / Thì đâu có cảnh huy hoàng ngày xuân.''\\[4pt]
\textit{(In the poem ``Encouraging Myself'' he wrote: `Without the bleakness of winter's end, / How could there be the brilliance of spring.')}

\end{truyen}
\begin{baihoc}
  Xà lim có thể giam giữ thân xác --- nhưng không thể giam giữ tinh thần và sức sáng tạo của người bất khuất.\\[4pt]
  \textit{(A cell can imprison the body --- but cannot imprison the spirit and creative power of the unyielding.)}
\end{baihoc}

\section{Helen Keller --- Vượt Qua Bóng Tối Và Im Lặng}
\begin{truyen}{Helen Keller --- Vượt Qua Bóng Tối Và Im Lặng}{Helen Keller --- Hoa Kỳ, 1880--1968}
\chuhoa{H}{elen} Keller bị mù và điếc hoàn toàn từ năm 19 tháng tuổi sau một trận sốt; cô sống trong thế giới tối đen và im lặng tuyệt đối suốt những năm đầu đời.\\[4pt]
\textit{(Helen Keller became completely blind and deaf at 19 months after a fever; she lived in a world of absolute darkness and silence through her early years.)}

Người thầy Anne Sullivan kiên nhẫn đánh vần từng chữ lên lòng bàn tay Helen, tuần này qua tuần khác, không bỏ cuộc.\\[4pt]
\textit{(Teacher Anne Sullivan patiently spelled each word into Helen's palm, week after week, never giving up.)}

Rồi một ngày --- ngày mà cô gọi là ngày thứ hai của cuộc đời --- khi nước chảy qua tay và Anne đánh vần ``W-A-T-E-R'', Helen bỗng hiểu rằng mọi thứ đều có tên.\\[4pt]
\textit{(Then one day --- the day she called the second day of her life --- water flowing over her hand as Anne spelled `W-A-T-E-R', Helen suddenly understood that everything has a name.)}

Helen Keller tốt nghiệp đại học, viết 12 cuốn sách, đi khắp thế giới truyền cảm hứng cho hàng triệu người.\\[4pt]
\textit{(Helen Keller graduated from university, wrote 12 books, traveled the world inspiring millions.)}

\end{truyen}
\begin{baihoc}
  Giới hạn của cơ thể không bao giờ là giới hạn của tinh thần --- nếu có một người thầy kiên nhẫn và một trái tim không chịu đầu hàng.\\[4pt]
  \textit{(The limits of the body are never the limits of the spirit --- when there is a patient teacher and a heart that refuses to surrender.)}
\end{baihoc}

\section{Mã Vân --- Từ Người Bị Từ Chối Đến Jack Ma}
\begin{truyen}{Mã Vân --- Từ Người Bị Từ Chối Đến Jack Ma}{Jack Ma --- Trung Quốc, 1964--}
\chuhoa{M}{ã} Vân thi trượt đại học ba lần, bị từ chối tất cả 31 vị trí xin việc kể cả nhân viên KFC, bị Harvard từ chối 10 lần, và thất bại trong hai công ty đầu tiên.\\[4pt]
\textit{(Ma Yun failed university entrance exams three times, was rejected for all 31 jobs he applied for including KFC, rejected by Harvard 10 times, and failed with his first two companies.)}

Khi thành lập Alibaba năm 1999 với 18 người trong căn hộ nhỏ, hầu hết mọi người đều coi đó là trò đùa.\\[4pt]
\textit{(When he founded Alibaba in 1999 with 18 people in a small apartment, most people considered it a joke.)}

Ông nói: ``Hôm nay thật khó khăn, ngày mai còn khó khăn hơn --- nhưng ngày kia sẽ thật tươi đẹp. Hầu hết mọi người từ bỏ vào ngày mai tối.''\\[4pt]
\textit{(He said: `Today is hard, tomorrow will be worse --- but the day after tomorrow will be beautiful. Most people give up on the dark tomorrow.')}

\end{truyen}
\begin{baihoc}
  Danh sách thất bại của người thành công thường dài hơn danh sách thành công của người thường --- sự khác biệt chỉ là họ không dừng lại.\\[4pt]
  \textit{(A successful person's list of failures is usually longer than an ordinary person's list of successes --- the only difference is they did not stop.)}
\end{baihoc}

\section{Phạm Tuân --- Từ Con Nông Dân Đến Du Hành Vũ Trụ}
\begin{truyen}{Phạm Tuân --- Từ Con Nông Dân Đến Du Hành Vũ Trụ}{Phạm Tuân --- Việt Nam, 1947--}
\chuhoa{P}{hạm} Tuân sinh ra trong gia đình nông dân nghèo ở Thái Bình, lớn lên trong thời chiến tranh, tình nguyện vào không quân năm 1965.\\[4pt]
\textit{(Pham Tuan was born into a poor farming family in Thai Binh, grew up during wartime, and volunteered for the air force in 1965.)}

Ngày 27 tháng 7 năm 1980, ông bay lên vũ trụ trên tàu Soyuz 37 --- trở thành người Châu Á đầu tiên bay vào vũ trụ.\\[4pt]
\textit{(On July 27, 1980, he launched into space aboard Soyuz 37 --- becoming the first Asian in space.)}

Từ những cánh đồng lúa Thái Bình đến vũ trụ bao la --- ông đã cho cả một thế hệ người Việt thấy rằng không có giấc mơ nào quá lớn khi đi kèm với ý chí đủ mạnh.\\[4pt]
\textit{(From the rice paddies of Thai Binh to the vast universe --- he showed an entire generation of Vietnamese that no dream is too large when paired with sufficient will.)}

\end{truyen}
\begin{baihoc}
  Xuất phát điểm không quyết định đích đến --- điều quyết định là bạn sẵn sàng đi bao xa và chịu đựng bao nhiêu để đến nơi mình muốn.\\[4pt]
  \textit{(Starting point does not determine destination --- what determines it is how far you are willing to go and how much you are willing to endure to get where you want.)}
\end{baihoc}
